\chapter{Obsah elektronické přílohy}
Součástí elektronické přílohy jsou primárně Matlab a Simulink soubory. Práci jsem zpracovával ve verzi Matlab R2025b, pro spuštění kódu proto doporučuji použít stejnou verzi. Skript param\_init není nutné manuálně spouštět, je volán automaticky při spuštění simulace.

\bigskip

{\small
%
\dirtree{%.
.1 /\DTcomment{kořenový adresář přiloženého archivu}.
.2 Horak.rizeni\_5L\_PMSM.pdf \DTcomment{elektronická verze závěrečné práce}.
.2 simulace.
.3 DP\_lib.slx \DTcomment{Knihovna se všemi vytvořenými funkcemi}.
.3 regulace\_rychlosti.slx \DTcomment{Simulink model celého řízení vč. simulace poskládaný z bloků pro generování HDL kódu}.
.3 param\_init.m \DTcomment{Skript inicializující parametry simulace}.
.4 Odvozene\_param.
.5 MTPA\_table.mat \DTcomment{Lookup tabulka MTPA}.
.5 sector\_LUT.mat.\DTcomment{Lookup tabulka pro výběr sektorů SVM modulátoru}.
.5 SVM\_const.mat.\DTcomment{Lookup tabulka napěťových vektorů pro SVM}.
.2 ovozeni\_a\_podpora.
.3 odvozeni.mlx.\DTcomment{Pomocné výpočty při odvozování}.
.3 mtpa\_init.m.\DTcomment{script generující MTPA tabulky}. 
.3 support\_svm.mlx.\DTcomment{script generující SVM lookup tabulku }.
}
}
